\chapter{拉取副本、ISR、高水位与确认}
本章从单机日志进入复制。\textbf{课程的分布式模型边界必须先读清：}三个独立 broker 各有 TCP server、线程与磁盘目录，follower 必须经 TCP 拉记录，不能共享内存复制；但它们运行在同一个 JVM，受一个可信 ClusterAuthority 管理。authority 线性化 metadata、epoch 与 HW 状态，保证教学切换可控。它不是 KRaft，不具备 controller 共识或 authority HA；整个 JVM/authority 丢失后的集群恢复不作保证。

复制只有在追加、报告和可见性边界彼此一致时才有意义。LEO 是每份日志自己的下一条 offset。ISR 是当前被允许参与确认的副本集合。HW 是普通消费者能读取的提交边界，消费者只能读取 offset < HW。副本通过真实网络拉取；authority 只提供一致状态和角色决策，不代替数据移动。

\begin{figure}[ht]
\centering
\begin{tikzpicture}[font=\small,node distance=20mm]
\node[draw,rounded corners,minimum width=31mm,minimum height=10mm] (leader) {leader log, LEO=5};
\node[draw,rounded corners,minimum width=31mm,minimum height=10mm,right=of leader] (follower) {follower log, LEO=3};
\node[draw,rounded corners,minimum width=33mm,minimum height=10mm,below=of $(leader)!0.5!(follower)$] (auth) {authority: ISR, HW=3};
\draw[-{Stealth},thick] (follower.north) to[bend left=18] node[above]{REPLICA\_FETCH} (leader.north);
\draw[-{Stealth},thick] (leader.south) -- node[left]{report LEO} (auth.north west);
\draw[-{Stealth},thick] (follower.south) -- node[right]{report LEO} (auth.north east);
\end{tikzpicture}
\caption{副本日志由 TCP 拉取，进度由可信 authority 串行化。示例 HW=min(5,3)=3。}
\end{figure}

\lessonsection{21}{拉取复制}
\goals{把“另一份日志”做成可验证的磁盘副本：follower 通过网络拉取 leader 记录并追加到自己的日志。}{理解副本为何有独立日志、复制读取为何能超过普通消费者的可见边界，以及一次 poll 不等于同步完成。}
\background{先修入口：课程步骤表按顺序要求完成\stepref{20}{消费组围栏}；实现时还要能读\stepref{6}{分段日志}的 \pathcode{PartitionLog} 与\stepref{12}{broker 请求}的 request/reply。leader 是当前接受分区写入的 broker；follower 是保存同一分区另一份日志、主动向 leader 拉记录的 broker。每台 broker 的日志独立落盘，复制不是共享对象或复制内存。LEO 是本地日志下一条记录的 offset；HW 是普通消费者可读的排他边界。follower 复制需要拿到 leader 的未提交尾部来追赶，所以 \pathcode{REPLICA_FETCH} 可读到 leader LEO，普通 consumer \pathcode{FETCH} 则只读 HW 以下。一次 poll 只完成“取回并本地追加”；第 22 步才把成功落盘的 LEO 报给 authority。}
\providededit{三 broker TCP harness、\pathcode{Messages.ReplicaFetchRequest} / \pathcode{ReplicaFetchBody}、\pathcode{ReplicaState}、\pathcode{PartitionLog} 生命周期及 framed RPC 已提供。fixture 的初始 authority reportedLEO 为 0；不能把直接写入 leader 磁盘误当作已报告同步。}{\pathcode{exercises/src/main/java/io/simplekafka/replication/FollowerReplicator.java}：\method{pollOnce(int maxRecords,int maxBytes)} 返回本次追加条数；并在 \pathcode{exercises/src/main/java/io/simplekafka/broker/BrokerHandler.java} 的 \method{handle(Messages.Request)} 中接通 \pathcode{REPLICA_FETCH} 分支。}
\paragraph{输入、输出与状态。}\pathcode{maxRecords} 是正的记录条数上限，\pathcode{maxBytes} 是正的编码字节预算，记录不能拆分；不足以容纳下一整条时可以暂时没有进展。构造时 broker id 非负并须与 \pathcode{ReplicaState} 的 id 相同；topic-partition、日志、借用的 RpcClient 与状态对象须非空。poll 返回实际追加数；空回复返回 0。RPC 客户端归调用者所有，poll 不关闭它。成功改变 follower 磁盘和本地 LEO；它不改 authority 的 reportedLEO、ISR 或 HW。
\lessoncontract{先在状态 monitor 下确认本地 broker 在线且为 follower，并捕获 epoch；在 monitor 外发请求；验证响应后再于同一 monitor 内复核 epoch/角色并追加。请求必须带当前分区、本地 LEO、broker id、捕获 epoch、两个正预算及 \pathcode{recoveryRead=false}。}
\begin{itemize}
\item 响应必须是副本读取 body，body epoch 与捕获 epoch 相同，leader LEO 不得低于本地 LEO，返回记录须从本地 LEO 开始、offset 连续且均小于 body 的 leader LEO。收到字节不代表已成功：只有 \pathcode{PartitionLog.append} 成功后才增加返回条数。
\item 非正预算或意外 body 为 \method{INVALID_REQUEST}；offline/已经成为 leader 为 \method{NOT_LEADER}；RPC 前后本地 epoch 改变或响应 epoch 错误为 \method{FENCED_EPOCH}；leader LEO 后退、缺 offset、错首 offset 或追加范围不对齐为 \method{CORRUPT_RECORD}。leader 端保留起点之外的远端错误（例如 \method{OFFSET_OUT_OF_RANGE}）及其 message 原样传播；本地存储错误保持原错误。
\item 任何上述验证失败均不得改 follower 日志；若网络请求期间角色变化，不得追加。leader 在 LEO 处返回空列表合法；反复空 poll 不改文件。broker 将过大的正字节预算压到单帧 reply 上限，后续 poll 继续复制剩余连续后缀。若本地写入发生 \method{STORAGE_ERROR}，不计入返回条数、也不报告进度；写入可能已有部分磁盘副作用，不能假定自动回滚，应遵循日志 failed/reopen 恢复路径。此步骤不调用 report，也不推进 HW。
\end{itemize}
\lessonhints{把“网络已返回”与“follower 磁盘追加成功”分开记录。}{逐项比对响应 epoch、首 offset、连续 offset 和 leader LEO；不要拿 HW 替代 LEO。}{对同一 follower 再 poll 一次，并分别造错 epoch、错首 offset、超出 leader LEO 的回复；检查 LEO 与文件字节是否保持不变。}
\expected{手算前提：分区起点为 0，epoch=0，broker 1 是 leader、broker 2 是 follower；leader 磁盘有三条普通记录，key 均为 null、value 分别是 1 个 ASCII 字节 \pathcode{a,b,c}、timestamp 为 11、12、13，故每条占 33 bytes；leader LEO=3、HW=0，authority 初始对三台 broker 的 reportedLEO 仍为 0。broker 2 调用 \pathcode{pollOnce(10,16384)} 收到 offset 0–2，追加三条后返回 3、本地 LEO=3；authority 不变。再次 poll 得空回复，返回 0 且文件字节不变。HW=0 时普通消费者仍看不到这些记录，但副本请求能读取它们。}
\paragraph{测试与完成标准。}真实 \method{Step21Test} 覆盖 \method{replicaFetchIncludesUncommittedLeaderTailAndReturnsEmptyAtLeo}、\method{followerCopySurvivesEmptyPollsAndCatalogReopen}、\method{boundsLargeReplicaFetchesAndCopiesTheCompleteSuffix}、\method{rejectedReplicaFetchRequestsPreserveDiskAndAuthority}、\method{followerAheadOfLeaderIsRejectedWithoutTruncatingItsLog} 与 \method{malformedRealTcpReplicaFetchRepliesAreRejectedBeforeAppend}：检查真实 TCP、磁盘字节、空 poll、单帧分页、follower领先时不截尾、epoch/offset 错误无副作用及重开内容。单步命令检查 Step21；累计命令从 Step1 跑到21。未实现骨架失败是预期，单步绿不表示第6章完成。
\pitfalls{直接共享 leader 的 \pathcode{PartitionLog} 绕过网络与 follower 落盘；把 HW 当成 leader LEO 会使 follower 无法复制尚未提交的尾部。真实 Kafka 的副本协议更复杂；这里的 \pathcode{REPLICA_FETCH} 是课程协议，不据此推断真实 Kafka 的内部帧或控制面。}
\lessoncommand{21}
\lessonquestions{为什么副本读取能取得 HW 以上、普通 consumer 暂时不可见的记录？}{为什么 follower 本地 append 成功之前不能把新 LEO 报给 authority？}
\paragraph{自检答案。}复制必须把 leader 未提交的尾部带到 follower，才能追赶当前 LEO；普通 consumer 的 HW 边界承担另一项工作。先持久化再报告可避免 authority 把尚未写入 follower 磁盘的进度计入确认。

\lessonsection{22}{ISR与高水位}
\goals{只用经过校验的进度报告计算消费者可见边界，并在受控时钟下移出长期未追上的 follower。}{理解物理 LEO、authority 报告、ISR、minISR 与 HW 不是同一个状态。}
\background{先修\stepref{21}{拉取复制}：follower 已在自己的日志写入记录，但该动作尚未报告。authority 的每分区状态保存各 broker 最近一次成功报告的 LEO；它与某份日志刚读出的物理 LEO 可以暂时不同。ISR 是当前参与确认的副本集合，minISR 是允许 ALL 请求继续的最少 ISR 数；HW 是 ISR 已共同确认前缀的排他末端。只有成功 report 才能改变 authority 进度。}
\providededit{每分区锁和 Condition、已分配 replica 集合、在线状态、authority 进度容器与可注入毫秒时钟 \pathcode{TimeSource} 已提供。}{\pathcode{exercises/src/main/java/io/simplekafka/replication/ReplicationTracker.java}：\method{report(int brokerId,int epoch,long leo)} 返回 void；\method{expireLagging()} 返回被移出 ISR 的 broker id 集合；\method{advanceHighWatermark()} 返回当前 HW。}
\paragraph{输入、输出与状态。}\pathcode{report} 的 broker 必须非负、已分配且在线，epoch/LEO 非负且 epoch 必须等于当前 leader epoch；无效输入不更改快照。leader 报告不得倒退；follower 报告不得超过 leader 已报告 LEO、也不得在同 epoch 倒退。成功报告更新该 replica 的 reportedLEO；follower 只有当 LEO 追到本次观察的 leader LEO 时才更新 last-caught-up 时间。报告本身不会把已移出 ISR 的成员重新加入。
\lessoncontract{若 ISR 含 leader、成员进度都存在且 ISR 大小至少 minISR，则候选 HW 是 ISR 中最新成功报告 LEO 的最小值，并应用 \pathcode{HW=max(oldHW,candidate)}；否则 HW 冻结。HW 单调不减。}
\begin{itemize}
\item 负 broker id、epoch 或 LEO，未知/未分配/离线 broker，leader LEO 倒退，或 follower 超前/倒退均为 \method{INVALID_REQUEST}；错误 epoch 为 \method{FENCED_EPOCH}。失败 report 不改变 reportedLEO、ISR 或 HW。
\item \pathcode{expireLagging()} 检查 ISR 中的非 leader；没有 catch-up 时间或 \pathcode{now-lastCaughtUp >= lagTimeoutMillis} 时移出 ISR，返回不可变的已移出集合。timeout=0 表示立即到期；now 早于 catch-up 时间时不按超时删除。过期不删除副本分配和日志，也不删除 leader。
\item 移出成员后重新考虑 HW；若余下 ISR 低于 minISR 则冻结。有效 report、实际过期或 HW 推进时通知等待者；\pathcode{advanceHighWatermark()} 按当前 ISR 重算并返回 HW，若 HW 实际推进会 signal 等待者；它本身不报告新的副本进度，也不执行 lag 过期。
\end{itemize}
\lessonhints{先验证 follower 报告的“不超过 leader”与“同 epoch 不倒退”，再计算 HW。}{minISR 不足时不要用已有副本的最小 LEO 推进边界；caught-up 时间只在追平时刷新。}{在固定 clock 上分别检查超时前一毫秒、恰好到期时刻及一个非法 report；比较快照而不靠 sleep。}
\expected{手算进度例：replicas/ISR={1,2,3}、minISR=2、旧 HW=0，leader 1 与 follower 2、3 的已报告 LEO 分别为5、3、4，则候选最小值为3，HW=3；follower 2 成功报告5后候选为4，HW=4；follower 3 再报告5后HW=5。单独看超时边界：follower 3 在clock=1000ms追平，lag timeout=1000ms，则now=1999ms尚保留，now=2000ms恰好到期并从ISR移除；不因它的LEO而改写日志。}
\paragraph{测试与完成标准。}\method{Step22Test.rejectsInvalidProgressAndExpiresOnlyReplicasThatMissTheExactCatchUpDeadline} 使用 clock=1000 起始，验证负进度、future epoch、未分配/离线、follower 超前及倒退均不改 snapshot；clock=1999 不过期、2000 移除 broker 3；也验证 ISR 缩小时 HW 冻结、重新报告推进到 4，而只剩 leader 时即使其报告到 5，HW 仍停在 4。单步命令只定位 Step 22；累计命令从 Step 1 到 22，单步通过不等于第 6 章完成。
\pitfalls{在线不等于 ISR，物理日志 LEO 不等于 authority 已知 LEO；用前者直接算 HW 会把未确认数据暴露给普通读取。课程用可信单 JVM authority 串行化状态；真实 Kafka 的跨 epoch HW 规则不保证本课程这里的全局单调性。}
\lessoncommand{22}
\lessonquestions{minISR 不足时为什么已有副本的 LEO 也不能推进 HW？}{某 follower 只是在线却没有追到 leader LEO，为什么不能刷新 catch-up 时间？}
\paragraph{自检答案。}HW 表示可读的共同前缀，不只是一个副本的写入位置；ISR 数不足时无法按本课程契约满足确认要求。只有实际追平的报告能证明 follower 已到当前 leader 尾部，在线连接本身不证明它持有该数据。

\lessonsection{23}{生产确认}
\goals{解释写入成功、复制提交与调用方收到确认的区别，并实现两种等待规则。}{理解 half-open offset 范围、Condition、单调时钟 deadline，以及超时不撤销 append。}
\background{先修\stepref{22}{ISR 与 HW 计算}。ACK 是 producer 要求 broker 等到哪一层再回复：\pathcode{LEADER} 表示 leader 本地追加完成即可返回；\pathcode{ALL} 表示需要 minISR 足够，且 HW 到达本批次的 \pathcode{nextOffset}。若返回范围为 [a,b)，需要 HW\(\ge b\) 才说明这批完整进入提交前缀；读者只能取 offset<b。deadline 是 \pathcode{System.nanoTime()} 时间域的绝对时刻，timeoutMillis 则是请求时长，不能互换。}
\providededit{ReplicationTracker 的 partition lock/Condition、\pathcode{AppendResult}、\pathcode{Acks}、\pathcode{RecordData} 与 \pathcode{TimeSource} 已提供；deadline helper \pathcode{deadlineAfterMillis(long)} 也已提供。}{\pathcode{exercises/src/main/java/io/simplekafka/replication/AckPolicy.java}：\method{validateBeforeAppend(Acks acks)} 返回 void；\method{await(AppendResult result,Acks acks,int epoch,long deadlineNanos)} 返回 void。}
\paragraph{输入、输出与副作用。}\pathcode{validateBeforeAppend} 接受非空的 LEADER/ALL；LEADER 不要求 ISR 检查，ALL 在追加前检查 ISR 大小。 \pathcode{await} 接受非空结果与 ACK 模式、该 append 捕获的 epoch、单调时钟绝对 deadline；LEADER 立即返回，ALL 等待相应 \pathcode{nextOffset}。调用端的 \pathcode{ReplicatedPartition.produce} 负责拒绝负 timeout 或 epoch。null 参数按接口抛 NullPointerException。任何追加之后的等待错误都不会撤回磁盘记录。
\lessoncontract{ALL 追加前 ISR<minISR 报 \method{NOT_ENOUGH_REPLICAS}，因此还没有追加；追加后按 \pathcode{epoch → ISR → HW → deadline} 重查：epoch 改变为 \method{FENCED_EPOCH}，ISR 低于 minISR 为 \method{NOT_ENOUGH_REPLICAS}，HW 已到结果 nextOffset 则成功，否则 deadline 到期为 \method{REQUEST_TIMEOUT}。HW 已满足时优先于过期 deadline 返回成功。Condition 等待释放 partition lock，醒来后必须重查；中断恢复线程 interrupt 标志并报 \method{REQUEST_TIMEOUT}，保留中断为 cause。LEADER 的 await 立即完成。}
\lessonhints{先分开“追加前 minISR 检查”和“追加后 HW 等待”，它们的副作用和错误不同。}{Condition 被 signal 只表示状态可能变化；重新取得锁后按优先顺序检查全部条件。}{用 latch 控制 report 顺序，验证 ALL 等待被报告唤醒；再用零 timeout、epoch 改变、ISR 下降和 interrupt 检查失败后的 LEO/HW。}
\expected{假设 leader epoch=0、ISR={1,2,3}、minISR=2，追加一条普通记录返回 [0,1)，其 key 为 null、value 为空、总长 32 bytes，leader LEO=1、follower LEO=0、HW=0。ALL timeout=0 在追加后报 \method{REQUEST_TIMEOUT}：LEO 保持 1、HW 保持 0、调用方没有 receipt，consumer 看不到 offset 0。稍后两个 follower 都复制并成功 report 到 1，HW 可推进到 1，记录变得可见；原请求不会追溯性地再得到成功 receipt。若在 append 前 ISR 已不足，ALL 则直接报 \method{NOT_ENOUGH_REPLICAS} 且 LEO 仍为 0。}
\paragraph{测试与完成标准。}\method{Step23Test} 包含 \method{insufficientIsrRejectsAllBeforeAppendWhileLeaderAckStillAppends}、\method{exactHighWatermarkBoundarySucceedsButOneOffsetBelowWaitsForProgress}、\method{zeroTimeoutRetainsTheAppendWithoutAdvancingTheHighWatermark}、\method{epochChangeWhileWaitingFencesTheAcknowledgement} 与 \method{isrReductionWhileWaitingFailsAndInterruptionPreservesTheInterruptFlag}，分别验证预检无副作用、HW 边界、超时保留数据、epoch/ISR 唤醒及中断标志。单步命令定位 Step 23；累计命令运行 Step 1–23，只有 Step 24 累计绿才完成本章。
\pitfalls{不能持 partition lock 等复制网络，否则复制报告永远无法取得锁；wall clock 会受校时跳变影响。timeout 只说明调用方没有按期限取得确认，并不说明 append 没发生；本课程仅提供两种 ACK 且不替调用方重试。}
\lessoncommand{23}
\lessonquestions{Condition 释放锁之后，为什么必须重查 epoch、ISR、HW 与 deadline？}{ALL 超时时，LEO、HW 与调用方知道写入结果之间是什么关系？}
\paragraph{自检答案。}等待期间任何状态都可能改变，signal 不是满足条件的证明。超时后 LEO 可已增加而 HW 未增加；因此数据可能稍后可见，但调用方仍没有可安全使用的成功 receipt。

\lessonsection{24}{并发副本broker}
\goals{把 leader 校验、写入、复制确认与普通读取组装成可并发服务的复制 broker。}{理解 authority lock 与本地日志锁顺序、HW 读取裁剪和 socket/Condition 生命周期。}
\background{先修\stepref{21}{拉取复制}、\stepref{22}{ISR 与 HW}及\stepref{23}{生产确认}。普通 consumer 只能读完整提交前缀 offset<HW；offset 在保留范围内但 \pathcode{offset >= HW} 时返回空列表。位于 LEO 之外的 offset 才越界。复制 fetch 是另一种请求，可为 follower 取未提交尾部。produce 的权限检查、append 与 leader progress 必须在 authority per-partition lock 保护下线性化；锁序是 authority lock→本地 PartitionLog monitor。网络调用不能持这两把锁；ALL 等待必须在锁外进行，让 follower report 能继续。}
\providededit{独立 TCP endpoints、broker worker、ReplicaState、ReplicationTracker、AckPolicy、PartitionLog 与 broker 路由已提供。}{\pathcode{exercises/src/main/java/io/simplekafka/replication/ReplicatedPartition.java}：\method{produce(List<RecordData> records,Acks acks,int epoch,long timeoutMillis)} 返回 \method{AppendResult}；\method{fetch(long offset,int maxRecords,int maxBytes,int epoch)} 返回 \pathcode{List<LogRecord>}。}
\paragraph{参数、返回与状态。}\pathcode{produce} 要求 records/acks 非空、epoch 和 timeoutMillis 非负、broker 当前在线且是请求 epoch 的 leader；成功返回 half-open append 范围，并可能追加日志、更新复制进度。 \pathcode{fetch} 要求 offset/epoch 非负、maxRecords/maxBytes 正；它只读取日志和 authority 状态，返回受记录数、编码字节预算与 HW 三重约束的完整记录，不拆单条记录。网络 worker 路由由提供的 handler 承担。
\lessoncontract{produce 先检查参数/null，再对 ALL 执行首次 append 前 minISR 检查；该检查发生在 online/角色/epoch 检查之前。随后在 authority lock 下同步并校验 online/leader/epoch，再次检查 ALL minISR（防止首次检查后状态变化），然后本地 append、leader report/可能推进 HW，最后解锁并等待 ACK。首次或第二次 minISR 检查失败、角色/epoch 失败均不得追加；首次检查可能使 minISR 错误先于 \method{NOT_LEADER}/\method{FENCED_EPOCH} 返回。负参数或空 produce batch 为 \method{INVALID_REQUEST}，非 leader 为 \method{NOT_LEADER}，旧 epoch 为 \method{FENCED_EPOCH}，存储错误为 \method{STORAGE_ERROR}。追加后 ACK 的 \method{NOT_ENOUGH_REPLICAS}、\method{REQUEST_TIMEOUT} 或 \method{FENCED_EPOCH} 不回滚 append。}
\begin{itemize}
\item fetch 对负 offset/epoch 或非正预算报 \method{INVALID_REQUEST}，对 stale epoch 报 \method{FENCED_EPOCH}，请求者不是在线 leader 报 \method{NOT_LEADER}；offset<logStart 或 offset>LEO 报 \method{OFFSET_OUT_OF_RANGE}。
\item fetch 的合法 offset 域是闭区间 [logStart,LEO]；域内且 offset\(\ge HW\) 返回空，offset<HW 时仍只返回 offset<HW 的整条记录。LEO 是下一 offset，故 offset=LEO 可以合法但为空；超过 LEO 不合法。
\item 锁序必须保持 authority→local log，不得倒序或持锁跨 broker RPC。关闭 broker 时必须让 socket/等待中的 ACK worker 结束；pending append 仍可能留在日志中。
\end{itemize}
\lessonhints{画出 produce、replica poll 与 consumer fetch 的锁范围，并把每次 TCP 调用放在图外。}{检查 HW 截断发生在普通 fetch 回复之前，而不是依赖 consumer 自己忽略数据。}{让 ALL producer 追加后暂停复制，同时从另一个连接 fetch，再并发写入 LEADER 记录；恢复复制后检查三者都能完成。}
\expected{并发例前提：leader=1、follower=2/3、epoch=0、初始空日志；先 LEADER 追加并复制 offset 0，三份 LEO=1、HW=1。再让 ALL 追加 offset 1，暂时 LEO leader=2、HW=1；另一个 fetch(0) 只能看到 offset 0，fetch(1) 为空。ALL 等待期间 LEADER 仍可追加 offset 2（范围 [2,3)），但 fetch 仍只见 offset 0。两个 follower 随后复制 offset 1–2 并报告 LEO=3，HW 推进到 3，ALL 返回 [1,2)，普通 fetch 才能看见 0、1、2。若每条普通记录 key=null、value 为空，则每条磁盘编码为 32 bytes：4 条中 HW=2 时，从 offset 0 且 maxBytes=31 返回空，预算32或63只容纳 offset0，预算64可容纳 offset0、1；maxRecords=1 即使预算更大也只返回一条。}
\paragraph{测试与完成标准。}\method{Step24Test} 的 \method{pendingAllAckLeavesFetchAtOldHwAllowsLeaderProduceAndCompletesAfterTcpReplication} 检查等待时旧 HW 可读与其他 produce 进展；\method{fetchHonorsHwLeoAndWholeRecordByteBudgetsAndRejectsInvalidRequests} 检查 31/32/63/64-byte 边界、HW/LEO/越界；\method{allPrecheckDoesNotAppendAndZeroTimeoutRetainsAnInvisibleAppendUntilReplication} 与 \method{stoppingLeaderCompletesAnOutstandingAllAckRequestWithoutLeakingTheProducer} 检查错误副作用和关闭唤醒。单步命令定位 Step 24；累计命令从 Step 1 到 24，是第 6 章完成边界。
\pitfalls{跨 socket 持锁会让唯一能推进复制的线程无法工作；HW 不是 group committed offset。此处三份独立日志经真实 TCP 复制，但 authority 仍是共享 JVM 内单点，因此不代表生产 HA。}
\lessoncommand{24}
\lessonquestions{ALL 仍等待时，为什么另一个 fetch 可安全读取旧 HW 前缀？}{关闭 broker 时只关闭 socket，为什么未必足以结束所有等待线程？}
\paragraph{自检答案。}fetch 读的是已经提交且不可变的旧前缀，未提交 offset 不会从回复泄漏。Condition waiter 还需被状态变化/关闭唤醒并检查退出条件；只关连接不能替代对等待线程和 worker 生命周期的处理。
